\section{Regression Machine Learning Model}
The advance of machine learning has led the wide applications in IoT devices.

Given the condition of running the models on constrained devices with
low computation power, memory and energy. The models must be
lightweight and fast during inference process, while also robust when
using interger instead
of floating point % TODO: talke about quanitisation
To be specific, the model must be compile in C with less than 10KB
and must not use floating point in its operation.
We have considered ... models namely: K-nearest neightbour
regression, linear regression, support vector regression, Gaussian
Process Regressor, Decision Tree, Random Forest and Multilayer perceptron model.

KNNR: works by comparing a node with the surrounding nodes to predict
a value. The closeness of values is typically calculated by Euclid
distance (...equation). Training in KNNR is very fast as the model
only remembers the training data. The inference is slow as the model
has to go through each point to calculate the distance and select the
best one. This means that KNNR does not fit our requirement of being
fast and lightweight as the model is heavy and computational heavy in inference.

However, KNNR is considered a lazy model as it has to remember all
the training data to calculate the distance with each node
remembers all training data, thus result in a heavy model and
contradict the low memory requirement.

Linear regression: works by fitting the data on a line, forming a
linear equation % TODO: insert latex equation for this
. This model satisfies the simplicity, requiring less computational
power and memory usage as the only storage is the coefficients,
however the linear regression model relies heavily on floating point
and is susceptible to category data type (hop count, neightbour
counts). Thus this is not a suitable model for the application

Likewise, SVM also falls short for similar reason as linear regression.

Gaussian Process Regressor, similar to SVM. However, in \cite{}, the
authors analyse GPR and conlcude that the memory requirement of GPR
is $O(n^2)$ while the computation cost is $O(n^3)$, which is quite
computational and memory expensive for constrained devices

Deep learning is recently a new adavnces and has achieve in terms of
matching the complex patterns and data, and works well for a wide
range of problems, including image, language and simple regression.
However, due to the constrained environment of the working MCU which
has several KB of RAM and ROM. I rule out this.

Random forest:
- Does not need scaling or normalisation
- Fairly simple
- Can work well with discrete value type:

Random forest and Decision Tree operates on similar algorithm. In
tree base algorithm, each value is split into left and right,
following the path to reach the end. In Random Forest, multiple trees
are used to prevent overfitting and provide better performance.
However, in my experiments, the performance of random forest and
decision tree are quite similar to each other. Thus I remove random
forest from list of considered models as they are much heavier than
decision tree.

In the end, we decide to proceed with Random Forest and
LGBMTreeRegressor due to the design and
computation simplicity to fit in with the constrained of the devices
that we operate on.
